% GENERATED FILE. Edit canonical lessons, then run npm run generate:books.
% canonical-source-hash: fnv1a64:627a2f4b46eae4a7
% canonical-lessons: MR-C362-andolan, MR-C362-nidarsan, MR-C362-uposan, MR-C362-nisedh-karne, MR-C362-bahiskar-ghalne

\chapter{Making Voices Heard}
\label{ch:mr-making-voices-heard}
\hlchaptermodality{362}

\begin{chapteropening}
\textbf{By the end of this chapter:} \emph{I can talk about movements, protests and boycotts.}

Complete the last lesson of chapter 362: I can talk about movements, protests and boycotts.
\end{chapteropening}

\section[\texorpdfstring{\mr{आंदोलन} (āndolan)}{आंदोलन (āndolan)}]{\mr{आंदोलन} (āndolan) — a movement; an agitation}

\label{lesson:MR-C362-andolan}



\begin{quote}
Before the new one: say the Marathi for a subsidy, then the Marathi for economic.
\end{quote}

\subsection*{You'll want to know: \mr{आंदोलन}}
\textbf{\mr{आंदोलन}} — \emph{āndolan} — \textquotedblleft{}a movement; an agitation\textquotedblright{}.

\begin{quote}
\mr{पाण्यासाठी} \mr{गावात} \mr{मोठं} \mr{आंदोलन} \mr{झालं}.
\end{quote}

\begin{quote}
\emph{pāṇyāsāṭhī gāvāt moṭhaṁ āndolan jhālaṁ.}
\end{quote}

\begin{quote}
There was a big agitation in the village for water.
\end{quote}

Neuter.

\mr{आंदोलन} \mr{करणे}, to agitate; \mr{आंदोलक}, an agitator.

\begin{cousinweb}
\textbf{\mr{आंदोलन}} (\emph{āndolan}) is a tatsama, taken from Sanskrit \emph{āndolana}, \textquotedblleft{}swinging, oscillation\textquotedblright{} (Monier-Williams). Marathi uses it for a movement that stirs people.
\end{cousinweb}

\begin{grammarlens}[title={What you've built}]
One more word for giving reasons and opinions.
\end{grammarlens}

\subsection*{Your turn}
\begin{itemize}
\raggedright
  \item \emph{Say it:} \emph{āndolan}
  \item \emph{Say it:} \emph{āndolan}, once more
  \item \emph{Say it:} \emph{pāṇyāsāṭhī gāvāt moṭhaṁ āndolan jhālaṁ.} — There was a big agitation in the village for water.
  \item \emph{Recall:} say the Marathi for a subsidy, then the Marathi for economic, then say \emph{āndolan} again
\end{itemize}

\begin{tcolorbox}[breakable,colback=teal!4,colframe=teal!35!black,title={Before you move on}]
What does \mr{आंदोलन} mean? (\textquotedblleft{}A movement; an agitation\textquotedblright{}.) Say it once more.
\end{tcolorbox}
\section[\texorpdfstring{\mr{निदर्शन} (nidarśan)}{निदर्शन (nidarśan)}]{\mr{निदर्शन} (nidarśan) — a demonstration (protest)}

\label{lesson:MR-C362-nidarsan}



\begin{quote}
Before the new one: say the Marathi for economic, then the Marathi for a movement.
\end{quote}

\subsection*{You'll want to know: \mr{निदर्शन}}
\textbf{\mr{निदर्शन}} — \emph{nidarśan} — \textquotedblleft{}a demonstration (protest)\textquotedblright{}.

\begin{quote}
\mr{कामगारांनी} \mr{कार्यालयासमोर} \mr{निदर्शनं} \mr{केली}.
\end{quote}

\begin{quote}
\emph{kāmgārānnī kāryālyāsamor nidarśanaṁ kelī.}
\end{quote}

\begin{quote}
The workers demonstrated in front of the office.
\end{quote}

Neuter. Plural: \mr{निदर्शनं}, \emph{nidarśanaṁ}.

Mostly plural: \mr{निदर्शनं} \mr{करणे}, to demonstrate.

\begin{cousinweb}
\textbf{\mr{निदर्शन}} (\emph{nidarśan}) is a tatsama, taken from Sanskrit \emph{nidarśana}, \textquotedblleft{}pointing to, showing\textquotedblright{} (Monier-Williams).
\end{cousinweb}

\begin{grammarlens}[title={What you've built}]
One more word for giving reasons and opinions.
\end{grammarlens}

\subsection*{Your turn}
\begin{itemize}
\raggedright
  \item \emph{Say it:} \emph{nidarśan}
  \item \emph{Say it:} \emph{nidarśan}, once more
  \item \emph{Say it:} \emph{kāmgārānnī kāryālyāsamor nidarśanaṁ kelī.} — The workers demonstrated in front of the office.
  \item \emph{Recall:} say the Marathi for economic, then the Marathi for a movement, then say \emph{nidarśan} again
\end{itemize}

\begin{tcolorbox}[breakable,colback=teal!4,colframe=teal!35!black,title={Before you move on}]
What does \mr{निदर्शन} mean? (\textquotedblleft{}A demonstration (protest)\textquotedblright{}.) Say it once more.
\end{tcolorbox}
\section[\texorpdfstring{\mr{उपोषण} (upoṣaṇ)}{उपोषण (upoṣaṇ)}]{\mr{उपोषण} (upoṣaṇ) — a hunger strike; a fast}

\label{lesson:MR-C362-uposan}



\begin{quote}
Before the new one: say the Marathi for a movement, then the Marathi for a demonstration.
\end{quote}

\subsection*{You'll want to know: \mr{उपोषण}}
\textbf{\mr{उपोषण}} — \emph{upoṣaṇ} — \textquotedblleft{}a hunger strike; a fast\textquotedblright{}.

\begin{quote}
\mr{मागण्यांसाठी} \mr{त्याने} \mr{तीन} \mr{दिवस} \mr{उपोषण} \mr{केलं}.
\end{quote}

\begin{quote}
\emph{māgaṇyāṁsāṭhī tyāne tīn divas upoṣaṇ kelaṁ.}
\end{quote}

\begin{quote}
He fasted for three days for the demands.
\end{quote}

Neuter.

\mr{उपोषण} \mr{करणे}, to go on a fast; \mr{उपास}, an ordinary fast.

\begin{cousinweb}
\textbf{\mr{उपोषण}} (\emph{upoṣaṇ}) is a tatsama, taken from Sanskrit \emph{upoṣaṇa}, \textquotedblleft{}fasting\textquotedblright{} (Monier-Williams).
\end{cousinweb}

\begin{grammarlens}[title={What you've built}]
One more word for giving reasons and opinions.
\end{grammarlens}

\subsection*{Your turn}
\begin{itemize}
\raggedright
  \item \emph{Say it:} \emph{upoṣaṇ}
  \item \emph{Say it:} \emph{upoṣaṇ}, once more
  \item \emph{Say it:} \emph{māgaṇyāṁsāṭhī tyāne tīn divas upoṣaṇ kelaṁ.} — He fasted for three days for the demands.
  \item \emph{Recall:} say the Marathi for a movement, then the Marathi for a demonstration, then say \emph{upoṣaṇ} again
\end{itemize}

\begin{tcolorbox}[breakable,colback=teal!4,colframe=teal!35!black,title={Before you move on}]
What does \mr{उपोषण} mean? (\textquotedblleft{}A hunger strike; a fast\textquotedblright{}.) Say it once more.
\end{tcolorbox}
\section[\texorpdfstring{\mr{निषेध} \mr{करणे} (niṣedh karṇe)}{निषेध करणे (niṣedh karṇe)}]{\mr{निषेध} \mr{करणे} (niṣedh karṇe) — to protest; to condemn}

\label{lesson:MR-C362-nisedh-karne}



\begin{quote}
Before the new one: say the Marathi for a demonstration, then the Marathi for a hunger strike.
\end{quote}

\subsection*{You'll want to know: \mr{निषेध} \mr{करणे}}
\textbf{\mr{निषेध} \mr{करणे}} — \emph{niṣedh karṇe} — \textquotedblleft{}to protest; to condemn\textquotedblright{}.

\begin{quote}
\mr{सगळ्यांनी} \mr{या} \mr{निर्णयाचा} \mr{निषेध} \mr{केला}.
\end{quote}

\begin{quote}
\emph{sagaḷyānnī yā nirṇayācā niṣedh kelā.}
\end{quote}

\begin{quote}
Everyone condemned this decision.
\end{quote}

Past, with -\mr{ने} (he …): \mr{त्याने} \mr{निषेध} \mr{केला}, \emph{tyāne niṣedh kelā}.

What is protested takes -\mr{चा}: \mr{निर्णयाचा} \mr{निषेध}.

\begin{cousinweb}
\textbf{\mr{निषेध} \mr{करणे}} (\emph{niṣedh karṇe}) joins the tatsama \mr{निषेध}, Sanskrit \emph{niṣedha}, \textquotedblleft{}keeping off, prohibition, denial\textquotedblright{} (Monier-Williams), to \mr{करणे}, \textquotedblleft{}to do\textquotedblright{}.
\end{cousinweb}

\begin{grammarlens}[title={What you've built}]
One more word for giving reasons and opinions.
\end{grammarlens}

\subsection*{Your turn}
\begin{itemize}
\raggedright
  \item \emph{Say it:} \emph{niṣedh karṇe}
  \item \emph{Say it:} \emph{niṣedh karṇe}, once more
  \item \emph{Say it:} \emph{sagaḷyānnī yā nirṇayācā niṣedh kelā.} — Everyone condemned this decision.
  \item \emph{Recall:} say the Marathi for a demonstration, then the Marathi for a hunger strike, then say \emph{niṣedh karṇe} again
\end{itemize}

\begin{tcolorbox}[breakable,colback=teal!4,colframe=teal!35!black,title={Before you move on}]
What does \mr{निषेध} \mr{करणे} mean? (\textquotedblleft{}To protest; to condemn\textquotedblright{}.) Say it once more.
\end{tcolorbox}
\section[\texorpdfstring{\mr{बहिष्कार} \mr{घालणे} (bahiṣkār ghālṇe)}{बहिष्कार घालणे (bahiṣkār ghālṇe)}]{\mr{बहिष्कार} \mr{घालणे} (bahiṣkār ghālṇe) — to boycott}

\label{lesson:MR-C362-bahiskar-ghalne}



\begin{quote}
Before the new one: say the Marathi for a hunger strike, then the Marathi for to protest.
\end{quote}

\subsection*{You'll want to know: \mr{बहिष्कार} \mr{घालणे}}
\textbf{\mr{बहिष्कार} \mr{घालणे}} — \emph{bahiṣkār ghālṇe} — \textquotedblleft{}to boycott\textquotedblright{}.

\begin{quote}
\mr{ग्राहकांनी} \mr{त्या} \mr{दुकानावर} \mr{बहिष्कार} \mr{घातला}.
\end{quote}

\begin{quote}
\emph{grāhkānnī tyā dukānāvar bahiṣkār ghātlā.}
\end{quote}

\begin{quote}
The customers boycotted that shop.
\end{quote}

Past, with -\mr{ने} (he …): \mr{त्याने} \mr{बहिष्कार} \mr{घातला}, \emph{tyāne bahiṣkār ghātlā}.

What is boycotted takes -\mr{वर}; \mr{घालणे}'s past is \mr{घातला}.

\begin{cousinweb}
\textbf{\mr{बहिष्कार} \mr{घालणे}} (\emph{bahiṣkār ghālṇe}) joins the tatsama \mr{बहिष्कार}, Sanskrit \emph{bahiṣkāra}, \textquotedblleft{}expulsion\textquotedblright{} (Monier-Williams), to \mr{घालणे}, \textquotedblleft{}to put on\textquotedblright{}.
\end{cousinweb}

\begin{grammarlens}[title={What you've built}]
One more word for giving reasons and opinions.
\end{grammarlens}

\subsection*{Your turn}
\begin{itemize}
\raggedright
  \item \emph{Say it:} \emph{bahiṣkār ghālṇe}
  \item \emph{Say it:} \emph{bahiṣkār ghālṇe}, once more
  \item \emph{Say it:} \emph{grāhkānnī tyā dukānāvar bahiṣkār ghātlā.} — The customers boycotted that shop.
  \item \emph{Recall:} say the Marathi for a hunger strike, then the Marathi for to protest, then say \emph{bahiṣkār ghālṇe} again
\end{itemize}

\begin{tcolorbox}[breakable,colback=teal!4,colframe=teal!35!black,title={Before you move on}]
What does \mr{बहिष्कार} \mr{घालणे} mean? (\textquotedblleft{}To boycott\textquotedblright{}.) Say it once more.
\end{tcolorbox}
